\SourceLeaf{438}%
\[
\begin{aligned}
\varrho d\xi&=Xds\\
\varrho d\eta&=Yds\\
\varrho d\zeta&=Zds
\end{aligned}
\]
and in place of $4$ curved lines on the surface of the auxiliary sphere, we have only $3$ to consider. Each element of the second line is thus to be regarded as lying in the great circle $L\lambda$, and the positive or negative value of $\varrho$ refers to the concavity or convexity of the curve in the direction of the normal.

\Needspace{5\baselineskip}
\begin{center}14.\end{center}
\noindent\hspace*{1.2em}We shall now investigate the spherical angle on the auxiliary sphere that the great circle from $L$ to $\lambda$ makes with the one from $L$ to one of the fixed points $(1),(2),(3)$, e.\,g. to $(3)$. To make this precise, we measure around from $L(3)$ to $L\lambda$ in the sense in which $(1),(2)$ and $(3)$ lie. If we call this angle $\varphi$, it follows by the proposition in Art. $7$ that
\[
\sin L(3).\sin L\lambda.\sin\varphi=Y\xi-X\eta
\]
or, since $L\lambda=90^{0}$ and
\[
\sin L(3)=\BareRoot{(XX+YY)}=\BareRoot{(1-ZZ)}:
\]
\[
\sin\varphi=\frac{Y\xi-X\eta}{\BareRoot{(XX+YY)}}.
\]
Furthermore,
\[
\sin L(3).\sin L\lambda.\cos\varphi=\zeta
\]
or
\[
\cos\varphi=\frac{\zeta}{\BareRoot{(XX+YY)}}
\]
and
\[
\operatorname{tang}\varphi=\frac{Y\xi-X\eta}{\zeta}=\frac{\zeta'}{\zeta}.
\]
Thus we have
\[
d\varphi=\frac{\zeta Yd\xi-\zeta Xd\eta-(Y\xi-X\eta)d\zeta+\xi\zeta dY-\eta\zeta dX}{(Y\xi-X\eta)^{2}+\zeta\zeta}.
\]
